\documentclass[11pt,aspectratio=169]{ctexbeamer}

\usetheme{default}
\setbeamertemplate{navigation symbols}{}

\usepackage{amsmath}
\usepackage{booktabs}

% Projector-safe palette: high contrast on common lecture-room projectors.
\definecolor{DeepBlue}{RGB}{11,58,102}
\definecolor{SoftInk}{RGB}{31,41,51}
\definecolor{PaperBg}{RGB}{255,255,255}
\definecolor{PanelBlue}{RGB}{226,237,247}

\setbeamercolor{normal text}{fg=SoftInk,bg=PaperBg}
\setbeamercolor{structure}{fg=DeepBlue}
\setbeamercolor{frametitle}{fg=DeepBlue,bg=PaperBg}
\setbeamercolor{title}{fg=DeepBlue}
\setbeamercolor{block title}{fg=DeepBlue,bg=PanelBlue}
\setbeamercolor{block body}{fg=SoftInk,bg=white}

\newcommand{\Core}[1]{
  \begin{block}{核心结论}
    #1
  \end{block}
}

\newcommand{\LectureMeta}[2]{
  \title{大模型推理系统与实践}
  \subtitle{#1}
  \author{张书豪}
  \institute{华中科技大学计算机学院}
  \date{2027 · #2}
}

\newcommand{\PaperRef}[3]{
  \item \textbf{#1}（#2）\hfill\href{#3}{原文链接}
}

\newcommand{\ResearchSources}[1]{
  \begin{frame}{本讲参考资料}
    \normalsize
    \begin{itemize}
      #1
    \end{itemize}
    \vspace{0.3em}
    \textcolor{SoftInk}{课件中的论文数字均指作者在特定模型、硬件和工作负载下报告的结果，不代表无条件性能。}
  \end{frame}
}

\newcommand{\InternalCaseNote}{
  \small\textcolor{SoftInk}{说明：本页为课程研究案例或教学假设，不代表已发表论文结论；只有在实验包给出模型、硬件、工作负载和原始记录时，数字才可作为结果引用。}
}

% 14 讲细分版的统一参考资料入口。论文数字的解释边界见课程内容基线。

\newcommand{\LectureOneSources}{\ResearchSources{
  \PaperRef{vLLM: Easy, Fast, and Cheap LLM Serving with PagedAttention}{SOSP 2023}{https://doi.org/10.1145/3600006.3613165}
  \PaperRef{SGLang: Efficient Execution of Structured Language Model Programs}{arXiv:2312.07104}{https://arxiv.org/abs/2312.07104}
  \PaperRef{nano-vllm-hust 教学代码}{课程仓库}{https://github.com/intellistream/intro-to-llm-inference-engines}
}}

\newcommand{\LectureTwoSources}{\ResearchSources{
  \PaperRef{vLLM / PagedAttention}{SOSP 2023}{https://doi.org/10.1145/3600006.3613165}
  \PaperRef{Sarathi-Serve}{OSDI 2024}{https://www.usenix.org/conference/osdi24/presentation/agrawal}
  \PaperRef{DistServe}{OSDI 2024}{https://www.usenix.org/conference/osdi24/presentation/zhong-yinmin}
}}

\newcommand{\LectureThreeSources}{\ResearchSources{
  \PaperRef{vLLM / PagedAttention}{SOSP 2023}{https://doi.org/10.1145/3600006.3613165}
  \PaperRef{DistServe}{OSDI 2024}{https://www.usenix.org/conference/osdi24/presentation/zhong-yinmin}
  \PaperRef{Llumnix}{OSDI 2024}{https://www.usenix.org/conference/osdi24/presentation/sun-biao}
}}

\newcommand{\LectureFourSources}{\ResearchSources{
  \PaperRef{Orca}{OSDI 2022}{https://www.usenix.org/conference/osdi22/presentation/yu}
  \PaperRef{Sarathi-Serve}{OSDI 2024}{https://www.usenix.org/conference/osdi24/presentation/agrawal}
  \PaperRef{FastServe}{NSDI 2026}{https://www.usenix.org/conference/nsdi26/presentation/wu-bingyang}
}}

\newcommand{\LectureFiveSources}{\ResearchSources{
  \PaperRef{vLLM / PagedAttention}{SOSP 2023}{https://doi.org/10.1145/3600006.3613165}
  \PaperRef{vAttention}{arXiv:2405.04437}{https://arxiv.org/abs/2405.04437}
  \PaperRef{Rethinking KV Cache Compression}{MLSys 2025}{https://proceedings.mlsys.org/paper_files/paper/2025/hash/26289c647c6828e862e271ca3c490486-Abstract-Conference.html}
  \PaperRef{Mooncake}{FAST 2025}{https://www.usenix.org/conference/fast25/presentation/qin}
}}

\newcommand{\LectureSixSources}{\ResearchSources{
  \PaperRef{DistServe}{OSDI 2024}{https://www.usenix.org/conference/osdi24/presentation/zhong-yinmin}
  \PaperRef{RAGCache}{arXiv:2404.12457}{https://arxiv.org/abs/2404.12457}
  \PaperRef{Mooncake}{FAST 2025}{https://www.usenix.org/conference/fast25/presentation/qin}
  \PaperRef{NEO}{MLSys 2025}{https://proceedings.mlsys.org/paper_files/paper/2025/hash/66a026c0d17040889b50f0dfa650e5e0-Abstract-Conference.html}
}}

\newcommand{\LectureSevenSources}{\ResearchSources{
  \PaperRef{vLLM / PagedAttention}{SOSP 2023}{https://doi.org/10.1145/3600006.3613165}
  \PaperRef{SGLang}{arXiv:2312.07104}{https://arxiv.org/abs/2312.07104}
  \PaperRef{DistServe}{OSDI 2024}{https://www.usenix.org/conference/osdi24/presentation/zhong-yinmin}
  \PaperRef{Mooncake}{FAST 2025}{https://www.usenix.org/conference/fast25/presentation/qin}
}}

\newcommand{\LectureEightSources}{\ResearchSources{
  \PaperRef{FlashAttention}{NeurIPS 2022}{https://arxiv.org/abs/2205.14135}
  \PaperRef{FlexGen}{ICML 2023}{https://proceedings.mlr.press/v202/sheng23a.html}
  \PaperRef{PowerInfer}{SOSP 2024 / arXiv:2312.12456}{https://arxiv.org/abs/2312.12456}
  \PaperRef{FlashInfer}{MLSys 2025}{https://proceedings.mlsys.org/paper_files/paper/2025/hash/dbf02b21d77409a2db30e56866a8ab3a-Abstract-Conference.html}
}}

\newcommand{\LectureNineSources}{\ResearchSources{
  \PaperRef{DEEPSERVE}{USENIX ATC 2025}{https://www.usenix.org/conference/atc25/presentation/hu-junhao}
  \PaperRef{QServe}{MLSys 2025}{https://proceedings.mlsys.org/paper_files/paper/2025/hash/fbe2b2f74a2ece8070d8fb073717bda6-Abstract-Conference.html}
  \PaperRef{vLLM-HUST}{开源维护仓库}{https://github.com/vLLM-HUST/vllm-hust}
}}

\newcommand{\LectureTenSources}{\ResearchSources{
  \PaperRef{vLLM / PagedAttention}{SOSP 2023}{https://doi.org/10.1145/3600006.3613165}
  \PaperRef{Revisiting Disaggregated LLM Serving}{2025 preprint, arXiv:2601.08833}{https://arxiv.org/abs/2601.08833}
  \PaperRef{Rethinking KV Cache Compression}{MLSys 2025}{https://proceedings.mlsys.org/paper_files/paper/2025/hash/26289c647c6828e862e271ca3c490486-Abstract-Conference.html}
}}

\newcommand{\LectureElevenSources}{\ResearchSources{
  \PaperRef{MLPerf Inference}{MLCommons benchmark}{https://mlcommons.org/benchmarks/inference-datacenter/}
  \PaperRef{ACM Artifact Review and Badging}{ACM policy}{https://www.acm.org/publications/policies/artifact-review-and-badging-current}
  \PaperRef{Revisiting Disaggregated LLM Serving}{2025 preprint, arXiv:2601.08833}{https://arxiv.org/abs/2601.08833}
}}

\newcommand{\LectureTwelveSources}{\ResearchSources{
  \PaperRef{nano-vllm-hust}{课程教学代码}{https://github.com/intellistream/intro-to-llm-inference-engines}
  \PaperRef{vLLM-HUST}{HUST 维护的 vLLM fork}{https://github.com/vLLM-HUST/vllm-hust}
  \PaperRef{上游 vLLM}{开源项目}{https://github.com/vllm-project/vllm}
}}

\newcommand{\LectureThirteenSources}{\ResearchSources{
  \PaperRef{课程内容基线与实验材料}{课程仓库}{https://github.com/intellistream/intro-to-llm-inference-engines}
  \PaperRef{DEEPSERVE}{USENIX ATC 2025}{https://www.usenix.org/conference/atc25/presentation/hu-junhao}
  \PaperRef{GhostServe}{MLSys 2026}{https://proceedings.mlsys.org/paper_files/paper/2026/hash/d99e8e80a6c41e148db686918dd7eab3-Abstract-Conference.html}
}}

\newcommand{\LectureFourteenSources}{\ResearchSources{
  \PaperRef{课程完整材料与后续更新}{课程仓库}{https://github.com/intellistream/intro-to-llm-inference-engines}
  \PaperRef{SpecGen}{MLSys 2026}{https://proceedings.mlsys.org/paper_files/paper/2026/hash/66a026c0d17040889b50f0dfa650e5e0-Abstract-Conference.html}
  \PaperRef{GhostServe}{MLSys 2026}{https://proceedings.mlsys.org/paper_files/paper/2026/hash/d99e8e80a6c41e148db686918dd7eab3-Abstract-Conference.html}
}}



\LectureMeta{第 6 讲 状态管理与记忆组织}{研究生课程课件}

\begin{document}

\begin{frame}[plain]
  \titlepage
\end{frame}

\begin{frame}{教学目标}
  \Core{本讲讨论推理系统中状态对象的保存、搬运、共享与复用问题。}
  \begin{itemize}
    \item 理解 DistServe 的阶段解耦。
    \item 理解 RAGCache 的知识对象缓存。
    \item 建立 KV、知识块、会话记忆的统一状态视角。
  \end{itemize}
\end{frame}

\begin{frame}{为什么状态问题会升级}
  \Core{当服务进入多轮交互、RAG 和智能体工作流后，状态不再是一次性中间结果。}
  \begin{itemize}
    \item 状态会跨请求、跨轮次持续存在。
    \item 状态对象可能被搬运到不同设备和不同阶段。
    \item 状态复用收益和一致性代价需要一起考虑。
  \end{itemize}
\end{frame}

\begin{frame}{状态对象有哪些类型}
  \Core{从系统角度看，不同状态对象的差别主要体现在生成方式、驻留时间和复用条件上。}
  \begin{itemize}
    \item 引擎内部状态：KV、激活、调度元数据。
    \item 外部知识状态：检索结果、知识块、索引缓存。
    \item 长程交互状态：会话历史、智能体记忆和工具上下文。
  \end{itemize}
\end{frame}

\begin{frame}{DistServe 的启发}
  \Core{DistServe 说明 prefill 和 decode 的分离可以带来收益，但状态传输会成为新的系统代价。}
  \begin{itemize}
    \item prefill 偏重算力，decode 偏重低排队时延和状态驻留。
    \item 解耦后可分别扩容和优化。
    \item 状态传输与协调复杂度随之出现。
  \end{itemize}
\end{frame}

\begin{frame}{DistServe 的系统启示}
  \Core{DistServe 的重要意义在于将阶段差异重新建模为资源配置与状态传递问题。}
  \begin{itemize}
    \item prefill 和 decode 不应强迫共享同一资源池。
    \item 一旦分离，状态搬运和阶段协同就进入主路径。
    \item 因此“解耦”本身不会免费带来收益。
  \end{itemize}
\end{frame}

\begin{frame}{RAGCache 的启发}
  \Core{RAGCache 说明外部知识一旦进入在线推理链路，就会变成系统内部必须管理的缓存对象。}
  \begin{itemize}
    \item 知识块注入会扩大 prefill 和 KV 压力。
    \item 热知识块存在跨请求复用机会。
    \item 检索与推理之间需要新的状态组织和重叠机制。
  \end{itemize}
\end{frame}

\begin{frame}{RAGCache 的系统启示}
  \Core{RAGCache 将知识注入过程转化为状态缓存与流水重叠问题。}
  \begin{itemize}
    \item 热知识块可以跨请求复用。
    \item 检索与推理需要联合优化而非串行拼接。
    \item 外部知识对象必须进入运行时视角。
  \end{itemize}
\end{frame}

\begin{frame}{统一状态对象视角}
  \Core{KV、知识缓存和会话记忆虽然服务不同语义目标，但在系统上都属于长期存在且可复用的状态对象。}
  \begin{itemize}
    \item 不同点在于复用条件、正确性要求和驻留时间。
    \item 相同点在于都需要考虑存储、调度、访问与回收。
    \item 统一视角有助于形成更强的系统抽象。
  \end{itemize}
\end{frame}

\begin{frame}{统一视角的研究价值}
  \Core{如果我们能把这些对象放进同一抽象框架，就有机会统一讨论存储、调度、迁移和复用。}
  \begin{itemize}
    \item 统一度量对象价值与代价。
    \item 统一考虑驻留层级和回收策略。
    \item 统一分析对调度器和执行路径的影响。
  \end{itemize}
\end{frame}

\begin{frame}{源码穿插：在 nano-vllm-hust 中观察状态对象}
  \Core{\texttt{\detokenize{nanovllm/engine/sequence.py}} 把请求状态、token 序列和 KV 相关元数据收在了同一个运行时对象里。}
  \begin{itemize}
    \item \texttt{\detokenize{Sequence}} 同时包含请求状态、token 序列和采样配置。
    \item \texttt{\detokenize{status}}、\texttt{\detokenize{num_cached_tokens}} 与 \texttt{\detokenize{block_table}} 体现了运行时状态的不同维度。
    \item \texttt{\detokenize{append_token()}}、\texttt{\detokenize{__getstate__()}} 与 \texttt{\detokenize{__setstate__()}} 说明状态对象会随阶段推进和跨进程序列化而变化。
  \end{itemize}
\end{frame}

\begin{frame}{源码穿插：Sequence 如何承载状态对象}
  \Core{通过 \texttt{\detokenize{Sequence}} 可以把“状态对象”从抽象概念落实到运行时最小承载单元。}
  \begin{itemize}
    \item 它既保存 prompt 与 completion 的边界，也保存调度阶段与缓存进度。
    \item \texttt{\detokenize{num_blocks}}、\texttt{\detokenize{last_block_num_tokens}} 等属性把逻辑请求与物理块管理关联起来。
    \item 这一设计有助于学生理解状态对象为何会同时受到调度、KV 管理和执行路径的共同影响。
  \end{itemize}
\end{frame}

\begin{frame}{讲解：统一状态视角到底统一了什么}
  \Core{这一讲不是把所有缓存混成一个词，而是把不同对象放到同一分析坐标里。}
  \begin{itemize}
    \item KV、检索知识块、会话记忆都具有跨时间存在、可能复用、需要回收这三个共同特征。
    \item 它们的差别在于生成方式、正确性要求、搬运代价和复用条件不同。
    \item 所谓“统一状态管理”，不是让它们实现完全一样，而是用同一套语言讨论驻留、迁移、命名和价值。
  \end{itemize}
\end{frame}

\begin{frame}{讲解：第六讲应形成的研究意识}
  \Core{状态一旦进入主路径，系统问题就不再只是计算快慢，而变成“对象如何持续存在”。}
  \begin{itemize}
    \item 只谈计算资源而不谈状态落点，很多机制会在真实系统里失效。
    \item 只谈潜在复用而不谈搬运与一致性，很多收益会停留在理想叙述里。
    \item 因此这一讲要让学生意识到：推理系统正在从执行系统走向状态化在线系统。
  \end{itemize}
\end{frame}

\begin{frame}{例子：为什么 DistServe 要把 prefill 和 decode 拆开（1）}
  \Core{状态管理不是“把对象存起来”，而是决定对象该在什么地方、什么时机被持有和搬运。}
\begin{block}{单机一体化路径}
    prefill 和 decode 共用一套资源，简单但容易让两类阶段互相干扰。
  \end{block}
\end{frame}

\begin{frame}{例子：为什么 DistServe 要把 prefill 和 decode 拆开（2）}
  \Core{状态管理不是“把对象存起来”，而是决定对象该在什么地方、什么时机被持有和搬运。}
\begin{block}{阶段解耦路径}
    prefill 更靠近高吞吐算力，decode 更靠近低抖动稳态资源，状态搬运因此成为主问题。
  \end{block}
\end{frame}

\begin{frame}{例子：为什么 DistServe 要把 prefill 和 decode 拆开（3）}
  \Core{状态管理不是“把对象存起来”，而是决定对象该在什么地方、什么时机被持有和搬运。}
\begin{block}{课堂结论}
    一旦阶段被拆开，系统研究的重点就从“存没存”转到“搬得动、对不对、值不值”。
  \end{block}
\end{frame}

\begin{frame}{例子：RAGCache 为什么让状态视角进一步扩张（1）}
  \Core{课程中的“状态”不只包括 KV，还包括检索知识块、会话记忆和调度元数据。}
\begin{block}{传统理解}
    缓存常被理解成底层实现优化，像一个被动的加速器。
  \end{block}
\end{frame}

\begin{frame}{例子：RAGCache 为什么让状态视角进一步扩张（2）}
  \Core{课程中的“状态”不只包括 KV，还包括检索知识块、会话记忆和调度元数据。}
\begin{block}{统一状态视角}
    检索结果、前缀块和会话历史都可能跨请求复用，因此它们都值得被当成一等状态对象建模。
  \end{block}
\end{frame}

\begin{frame}{例子：RAGCache 为什么让状态视角进一步扩张（3）}
  \Core{课程中的“状态”不只包括 KV，还包括检索知识块、会话记忆和调度元数据。}
\begin{block}{系统含义}
    一旦对象具有跨时间和跨请求价值，系统就必须管理它的命名、引用、迁移和回收边界。
  \end{block}
\end{frame}

\begin{frame}{代码例子：Sequence 的序列化为什么值得课堂关注（1）}
  \Core{状态管理课程里的代码例子，应让学生看到对象不只是运行时临时变量，还需要可传递与可恢复。}
\begin{block}{代码片段}
    {\ttfamily\small
    \detokenize{def __getstate__(self): ...}\par
    \detokenize{def __setstate__(self, state): ...}\par
    }
  \end{block}
\end{frame}

\begin{frame}{代码例子：Sequence 的序列化为什么值得课堂关注（2）}
  \Core{状态管理课程里的代码例子，应让学生看到对象不只是运行时临时变量，还需要可传递与可恢复。}
\begin{block}{课堂应追问}
    哪些字段如果丢失，请求将无法继续？哪些字段如果错传，会直接造成状态不一致？
  \end{block}
\end{frame}

\begin{frame}{代码例子：Sequence 的序列化为什么值得课堂关注（3）}
  \Core{状态管理课程里的代码例子，应让学生看到对象不只是运行时临时变量，还需要可传递与可恢复。}
\begin{block}{系统含义}
    只要对象可能跨模块、跨节点或跨阶段流动，序列化语义就是系统正确性的一部分。
  \end{block}
\end{frame}

\begin{frame}{代码例子：统一状态管理器会暴露哪些接口（1）}
  \Core{课程项目里一个常见扩展，是想象如果把 KV、知识块和会话记忆统一管理，需要哪些最小接口。}
\begin{block}{接口草图}
    {\ttfamily\small
    \detokenize{state_id = state_store.put(obj, metadata)}\par
    \detokenize{state_store.pin(state_id, device="gpu")}\par
    \detokenize{state_store.evict(policy="lru") }\par
    }
  \end{block}
\end{frame}

\begin{frame}{代码例子：统一状态管理器会暴露哪些接口（2）}
  \Core{课程项目里一个常见扩展，是想象如果把 KV、知识块和会话记忆统一管理，需要哪些最小接口。}
\begin{block}{课堂应追问}
    这些接口背后分别对应命名、驻留、迁移还是回收问题？
  \end{block}
\end{frame}

\begin{frame}{代码例子：统一状态管理器会暴露哪些接口（3）}
  \Core{课程项目里一个常见扩展，是想象如果把 KV、知识块和会话记忆统一管理，需要哪些最小接口。}
\begin{block}{结论}
    统一状态管理的难点不是接口名字，而是不同状态对象的价值函数和一致性需求并不相同。
  \end{block}
\end{frame}

\begin{frame}{例子：状态搬运为什么经常吞掉理论收益（1）}
  \Core{很多机制论文的理论收益很强，但一旦进入真实系统，搬运与同步代价就会迅速显形。}
\begin{block}{理想世界}
    把更适合 prefill 的资源和更适合 decode 的资源拆开，看起来可以两头占优。
  \end{block}
\end{frame}

\begin{frame}{例子：状态搬运为什么经常吞掉理论收益（2）}
  \Core{很多机制论文的理论收益很强，但一旦进入真实系统，搬运与同步代价就会迅速显形。}
\begin{block}{真实世界}
    KV 或知识对象一旦要跨节点移动，带宽、延迟、重组和一致性全都会冒出来。
  \end{block}
\end{frame}

\begin{frame}{例子：状态搬运为什么经常吞掉理论收益（3）}
  \Core{很多机制论文的理论收益很强，但一旦进入真实系统，搬运与同步代价就会迅速显形。}
\begin{block}{教学重点}
    讲状态复用时必须同时讲状态搬运，否则学生会高估机制落地的直接性。
  \end{block}
\end{frame}

\begin{frame}{课堂练习：把三类状态放到同一张表里（1）}
  \Core{学生需要学会比较不同状态对象，而不是只会孤立讨论 KV。}
\begin{block}{练习一}
    比较 KV、RAG 检索块、会话记忆三者在生成方式、驻留时间和复用条件上的差异。
  \end{block}
\end{frame}

\begin{frame}{课堂练习：把三类状态放到同一张表里（2）}
  \Core{学生需要学会比较不同状态对象，而不是只会孤立讨论 KV。}
\begin{block}{练习二}
    说明哪一类状态最适合先做 observe-only 探针，哪一类最需要严格正确性边界。
  \end{block}
\end{frame}

\begin{frame}{课堂练习：把三类状态放到同一张表里（3）}
  \Core{学生需要学会比较不同状态对象，而不是只会孤立讨论 KV。}
\begin{block}{练习三}
    用一句话回答：为什么第 6 讲的关键词不是“缓存”，而是“状态组织”？
  \end{block}
\end{frame}


\begin{frame}{课堂讨论}
  \Core{状态管理的关键难点，在于缓存对象如何稳定纳入系统控制环。}
  \begin{itemize}
    \item 为什么知识缓存不能被简单看成外挂组件？
    \item 当状态搬运进入主路径后，如何权衡收益与代价？
    \item 会话记忆在系统层面与 KV 有哪些相同点和差异？
  \end{itemize}
\end{frame}

\begin{frame}{本讲总结}
  \Core{DistServe 与 RAGCache 都要求系统显式管理状态，但 KV、检索文档和长期记忆的复用键、失效条件与安全边界并不相同。}
  \begin{itemize}
    \item 状态管理是推理基础设施的核心问题之一。
    \item 后续架构与异构执行讨论都要继续围绕状态展开。
  \end{itemize}
\end{frame}

\LectureSixSources

\begin{frame}{对应 Tutorial}
  \Core{本讲对应 Tutorial 5。先独立作答，再使用对应 Tutorial 的教师讨论要点核对。}
  \begin{itemize}
    \item 文件：\texttt{\detokenize{build/tutorials/Tutorial_05_状态管理与记忆组织.pdf}}
    \item 动手运行、日志记录与提交要求统一查看 \texttt{\detokenize{build/experiments/}} 和对应 experiment 源目录下的 \texttt{\detokenize{assignment_spring-2027/}}。
  \end{itemize}
\end{frame}

\end{document}
